% ============================================================
%  Section 1 — Introduction
% ============================================================
\section{Introduction}
\label{sec:intro}

An accountant processing a stack of Tunisian invoices faces the same four steps
on every document: read the supplier name, extract the amounts, find the right
account code from the company's chart, and check whether the numbers add up.
Repeated hundreds of times a week, this costs 3–5 minutes per document and
introduces a steady stream of classification errors.

Large Language Models (LLMs) can now perform all four steps in seconds. The challenge
is not capability — it is \textit{control}. A production system must:

\begin{itemize}[leftmargin=*, itemsep=2pt]
  \item assign codes from the \emph{company's own} chart, not a generic one;
  \item provide an auditable record showing which rule was used and why;
  \item flag documents with arithmetic errors or missing tax identifiers before
        they enter the ledger;
  \item update its behaviour when accounting rules change, without retraining.
\end{itemize}

DocDesk addresses all four requirements. It orchestrates five specialized agents
through n8n~\cite{n8n2024}, grounds account code assignment in a company-specific
rules store via Retrieval-Augmented Generation (RAG)~\cite{lewis2020rag}, and
enforces arithmetic and identifier checks in deterministic Python code.
Company-specific behaviour is stored in versioned skills. SkillOpt~\cite{skillopt2024}
is integrated as a text-space optimizer. Its failure-only update policy keeps
candidate edits bounded and accepts them only when the held-out selection score
improves. We also evaluate a separate, manually authored TRAIN-derived skill
revision, clearly distinguished from the SkillOpt run.

\textbf{Why not a single LLM call?}
Routing, extraction, coding, and validation are different tasks with different
failure modes. Separating them allows stage-level outputs, checks, and traces,
but adds orchestration complexity and can increase latency. In this workflow,
routing is deterministic and validation calls an LLM only when a check fails;
the pipeline does not make five model calls for every document.

\textbf{Contributions.}
\begin{enumerate}[leftmargin=*, itemsep=2pt]
  \item An n8n-based prototype connecting multimodal document intake,
        company-filtered rule retrieval, validation, review, and audit logging
        (\cref{sec:architecture}).
  \item A 300-document synthetic corpus with injected errors and hidden house
        rules designed to stress-test every agent (\cref{sec:data}).
  \item Quantitative Coder and Validator baselines, with a small observed
        with/without-RAG comparison; the broad Extractor metric is withheld
        pending a corrected rerun (\cref{sec:evaluation}).
  \item A measured baseline, prompt iteration, and a real SkillOpt run with a
        selection-gated no-change outcome on a balanced pilot split
        (\cref{sec:evaluation,sec:skillopt}).
\end{enumerate}
